\SourceLeaf{378}%
\needspace{8\baselineskip}
\begin{center}[2.]\end{center}

The general solution to the problem of mapping a curved surface onto a plane with similarity in the smallest parts appears to be as follows.

Let $V=0$ be the equation of the curved surface and let
\[
dV=p\,dx+q\,dy+r\,dz.
\]
Integrate the equation
\[
q\,dx-p\,dy+\BareRoot{-1}\,.\,\BareRoot{(pp+qq+rr)}\,.\,dz=0,
\]
by combining it with $V=0$. Let the integral be
\[
W=\mathrm{Const.}
\]
\Needspace{9\baselineskip}
Similarly, take the integral of
\[
q\,dx-p\,dy-\BareRoot{-1}\,.\,\BareRoot{(pp+qq+rr)}\,.\,dz=0:
\]
to be
\[
W'=\mathrm{Const.}
\]
The complete solution is then
\[
\begin{aligned}
\xi+\BareRoot{-1}\,.\,\eta&=f(W)\\
\xi-\BareRoot{-1}\,.\,\eta&=f'(W),
\end{aligned}
\]
where $f,f'$ denote arbitrary functions. Or rather, let the integration of
\[
q\,dx-p\,dy+\BareRoot{-1}\,.\,\BareRoot{(pp+qq+rr)}\,.\,dz=0
\]
give
\[
T+U\BareRoot{-1}=\mathrm{Const.},
\]
Then
\[
\xi\pm\BareRoot{-1}\,.\,\eta=\mathrm{funct.}\,(T\pm U\BareRoot{-1}).
\]

If the integral of
\[
dx^2+dy^2+dz^2=0
\]
in combination with $V=0$ is
\[
T+U\BareRoot{-1}=\mathrm{Const.},
\]
then the result above holds.
